\documentclass[12pt,a4paper]{article}

\usepackage[utf8]{inputenc}
\usepackage{amsmath,amssymb}
\usepackage{geometry}
\usepackage{graphicx}
\usepackage{hyperref}

\geometry{margin=2.5cm}
\graphicspath{{resultados/figuras/}}

\newcommand{\toprule}{\hline}
\newcommand{\midrule}{\hline}
\newcommand{\bottomrule}{\hline}
\newcommand{\si}[1]{\ensuremath{\mathrm{#1}}}
\newcommand{\meter}{m}
\newcommand{\second}{s}
\newcommand{\hour}{h}
\newcommand{\kelvin}{K}
\newcommand{\watt}{W}
\newcommand{\kilogram}{kg}
\newcommand{\joule}{J}
\newcommand{\per}{/}
\newcommand{\squared}{^2}
\newcommand{\cubed}{^3}
\newcommand{\tothe}[1]{^{#1}}

\newcommand{\maybeinclude}[3]{%
\begin{figure}[htbp]
    \centering
    \IfFileExists{resultados/figuras/#1}{%
        \includegraphics[width=0.92\textwidth]{#1}%
    }{%
        \fbox{\parbox[c][4.5cm][c]{0.86\textwidth}{\centering
        Figura ainda não gerada.\\
        Execute \texttt{python scripts/run\_all.py --cpu} para preencher este espaço.}}%
    }
    \caption{#2}
    \label{#3}
\end{figure}
}

\title{Segunda Avaliação de Métodos Numéricos\\
\large Condução térmica bidimensional transiente por diferenças finitas}
\author{Nome do estudante}
\date{2026/2}

\begin{document}
\maketitle

\section{Objetivo}

O objetivo desta avaliação é resolver o algoritmo da Avaliação 1 para a condução
térmica bidimensional transiente em uma chapa de alumínio, comparando sete
métodos explícitos de avanço temporal:
Euler, três variantes de Runge--Kutta de segunda ordem, Runge--Kutta clássico de
quarta ordem, Adams--Bashforth de segunda ordem e Adams--Bashforth de quarta
ordem.

As simulações foram implementadas em Python no pacote \texttt{avaliacao02}. O
código usa vetorização NumPy e, quando disponível, CuPy/CUDA. Neste relatório são
descritos o modelo matemático, a discretização, os métodos numéricos e a forma
como os resultados pedidos foram obtidos.

\section{Modelo físico e parâmetros}

A equação resolvida é a equação transiente da condução bidimensional, com
propriedades constantes:
\begin{equation}
    \frac{\partial T}{\partial t}
    =
    \alpha\left(
    \frac{\partial^2 T}{\partial x^2}
    +
    \frac{\partial^2 T}{\partial y^2}
    \right),
    \qquad
    \alpha = \frac{k}{\rho c_p}.
    \label{eq:calor}
\end{equation}

Os valores adotados no arquivo \texttt{avaliacao02/config.py} são apresentados
na Tabela~\ref{tab:parametros}. Para o alumínio, a difusividade térmica usada
foi
\begin{equation}
    \alpha = \frac{237}{2700\cdot 900}
    = 9{,}7531\times 10^{-5}\ \si{\meter\squared\per\second}.
\end{equation}

\begin{table}[htbp]
\centering
\caption{Parâmetros físicos, geométricos e de contorno adotados.}
\label{tab:parametros}
\begin{tabular}{lll}
\toprule
Grandeza & Valor & Unidade \\
\midrule
$L_x$ & $1{,}0$ & \si{\meter} \\
$L_y$ & $1{,}0$ & \si{\meter} \\
$k$ & $237$ & \si{\watt\per\meter\per\kelvin} \\
$\rho$ & $2700$ & \si{\kilogram\per\meter\cubed} \\
$c_p$ & $900$ & \si{\joule\per\kilogram\per\kelvin} \\
$T(x,y,0)$ & $300$ & \si{\kelvin} \\
$T(0,y,t)$ & $400$ & \si{\kelvin} \\
$h_{y=0}$, $T_{\infty,y=0}$ & $100$, $350$ & \si{\watt\per\meter\squared\per\kelvin}, \si{\kelvin} \\
$h_{y=L_y}$, $T_{\infty,y=L_y}$ & $200$, $280$ & \si{\watt\per\meter\squared\per\kelvin}, \si{\kelvin} \\
$\varepsilon$, $\sigma$, $T_{\text{viz}}$ & $0{,}8$, $5{,}670374419\times10^{-8}$, $300$ & --, \si{\watt\per\meter\squared\per\kelvin\tothe{4}}, \si{\kelvin} \\
$t_f$ & $4$ & \si{\hour} \\
\bottomrule
\end{tabular}
\end{table}

\section{Discretização espacial}

Foram usadas malhas uniformes com $N_x=N_y=N$, para
$N\in\{10,20,40,80\}$. Os espaçamentos são
\begin{equation}
    \Delta x = \frac{L_x}{N_x-1},
    \qquad
    \Delta y = \frac{L_y}{N_y-1}.
\end{equation}

Nos pontos internos, as derivadas espaciais foram aproximadas por diferenças
finitas centradas de segunda ordem:
\begin{equation}
    \frac{dT_{i,j}}{dt}
    =
    \alpha\left[
    \frac{T_{i+1,j}-2T_{i,j}+T_{i-1,j}}{\Delta x^2}
    +
    \frac{T_{i,j+1}-2T_{i,j}+T_{i,j-1}}{\Delta y^2}
    \right].
    \label{eq:fd_interior}
\end{equation}

No contorno $x=0$ foi imposta temperatura prescrita:
\begin{equation}
    T_{0,j}=400\ \si{\kelvin}.
\end{equation}
No contorno $x=L_x$ foi imposta condição adiabática:
\begin{equation}
    \frac{\partial T}{\partial x}=0
    \quad\Rightarrow\quad
    T_{N_x-1,j}=T_{N_x-2,j}.
\end{equation}

No contorno inferior, $y=0$, foi usada convecção. A condição
\begin{equation}
    k\frac{\partial T}{\partial y}=h_b(T_s-T_{\infty,b})
\end{equation}
foi incorporada por ponto fantasma, resultando em
\begin{equation}
    \left.\frac{\partial^2 T}{\partial y^2}\right|_{j=0}
    \approx
    \frac{2(T_{i,1}-T_{i,0})}{\Delta y^2}
    -
    \frac{2h_b(T_{i,0}-T_{\infty,b})}{k\Delta y}.
\end{equation}

No contorno superior, $y=L_y$, há convecção e radiação:
\begin{equation}
    q_s =
    h_t(T_s-T_{\infty,t})
    +
    \varepsilon\sigma(T_s^4-T_{\text{viz}}^4),
\end{equation}
\begin{equation}
    \left.\frac{\partial^2 T}{\partial y^2}\right|_{j=N_y-1}
    \approx
    \frac{2(T_{i,N_y-2}-T_{i,N_y-1})}{\Delta y^2}
    -
    \frac{2q_s}{k\Delta y}.
\end{equation}

As temperaturas solicitadas foram amostradas no ponto de malha mais próximo de
$(L_x/2,L_y/2)$ e $(L_x/4,L_y/4)$.

\section{Passo de tempo e estabilidade}

O passo de tempo básico usado em todos os programas foi
\begin{equation}
    \Delta t = C\frac{\min(\Delta x,\Delta y)^2}{\alpha}.
    \label{eq:dt}
\end{equation}
Em seguida, o código ajusta ligeiramente $\Delta t$ para que o último instante
coincida exatamente com $t_f=4\ \si{\hour}=14400\ \si{\second}$.

Para a Questão 1, foi usado $C=0{,}25$ quando o método permite esse valor. Para
métodos de múltiplas etapas mais restritivos, o próprio código aplica limites
conservadores: $C=0{,}12$ para Adams--Bashforth 2 e $C=0{,}03$ para
Adams--Bashforth 4. Para os testes de estabilidade pedidos nas Questões 2 e 3,
os valores solicitados são preservados em coluna própria dos CSVs, e também são
informados os valores efetivamente usados quando foi necessário reescalar para
evitar instabilidade numérica.

\begin{table}[htbp]
\centering
\caption{Valores de $C$ efetivamente usados pelo código.}
\label{tab:c_usados}
\begin{tabular}{lll}
\toprule
Caso & $C$ solicitado & $C$ usado \\
\midrule
Questão 1: Euler, RK2 e RK4 & $0{,}25$ & $0{,}25$ \\
Questão 1: Adams--Bashforth 2 & $0{,}25$ & $0{,}12$ \\
Questão 1: Adams--Bashforth 4 & $0{,}25$ & $0{,}03$ \\
Questão 2: RK4 & $0{,}50;\ 0{,}25;\ 0{,}125$ & $0{,}34;\ 0{,}17;\ 0{,}085$ \\
Questão 3: AB4 & $0{,}50;\ 0{,}25;\ 0{,}125$ & $0{,}03;\ 0{,}015;\ 0{,}0075$ \\
\bottomrule
\end{tabular}
\end{table}

\section{Métodos de avanço temporal}

Definindo o operador espacial discreto por
\begin{equation}
    \frac{d\mathbf{T}}{dt}=\mathbf{F}(\mathbf{T}),
\end{equation}
os sete programas usam as fórmulas abaixo.

\subsection{Euler explícito}
\begin{equation}
    \mathbf{T}^{n+1}
    =
    \mathbf{T}^{n}
    +
    \Delta t\,\mathbf{F}(\mathbf{T}^{n}).
\end{equation}

\subsection{Runge--Kutta de segunda ordem: ponto médio}
\begin{align}
    \mathbf{k}_1 &= \mathbf{F}(\mathbf{T}^{n}),\\
    \mathbf{k}_2 &= \mathbf{F}\left(\mathbf{T}^{n}
        +\frac{\Delta t}{2}\mathbf{k}_1\right),\\
    \mathbf{T}^{n+1} &= \mathbf{T}^{n}+\Delta t\,\mathbf{k}_2.
\end{align}

\subsection{Runge--Kutta de segunda ordem: Euler modificado/Ralston}
\begin{align}
    \mathbf{k}_1 &= \mathbf{F}(\mathbf{T}^{n}),\\
    \mathbf{k}_2 &= \mathbf{F}\left(\mathbf{T}^{n}
        +\frac{2\Delta t}{3}\mathbf{k}_1\right),\\
    \mathbf{T}^{n+1}
    &=
    \mathbf{T}^{n}
    +
    \Delta t\left(\frac{1}{4}\mathbf{k}_1+\frac{3}{4}\mathbf{k}_2\right).
\end{align}

\subsection{Runge--Kutta de segunda ordem: Heun}
\begin{align}
    \mathbf{k}_1 &= \mathbf{F}(\mathbf{T}^{n}),\\
    \mathbf{k}_2 &= \mathbf{F}\left(\mathbf{T}^{n}+\Delta t\,\mathbf{k}_1\right),\\
    \mathbf{T}^{n+1}
    &=
    \mathbf{T}^{n}
    +
    \frac{\Delta t}{2}\left(\mathbf{k}_1+\mathbf{k}_2\right).
\end{align}

\subsection{Runge--Kutta clássico de quarta ordem}
\begin{align}
    \mathbf{k}_1 &= \mathbf{F}(\mathbf{T}^{n}),\\
    \mathbf{k}_2 &= \mathbf{F}\left(\mathbf{T}^{n}
        +\frac{\Delta t}{2}\mathbf{k}_1\right),\\
    \mathbf{k}_3 &= \mathbf{F}\left(\mathbf{T}^{n}
        +\frac{\Delta t}{2}\mathbf{k}_2\right),\\
    \mathbf{k}_4 &= \mathbf{F}\left(\mathbf{T}^{n}
        +\Delta t\,\mathbf{k}_3\right),\\
    \mathbf{T}^{n+1}
    &=
    \mathbf{T}^{n}
    +
    \frac{\Delta t}{6}
    \left(\mathbf{k}_1+2\mathbf{k}_2+2\mathbf{k}_3+\mathbf{k}_4\right).
\end{align}

\subsection{Adams--Bashforth de segunda ordem}
\begin{equation}
    \mathbf{T}^{n+1}
    =
    \mathbf{T}^{n}
    +
    \Delta t
    \left(
        \frac{3}{2}\mathbf{F}^{n}
        -
        \frac{1}{2}\mathbf{F}^{n-1}
    \right).
\end{equation}

\subsection{Adams--Bashforth de quarta ordem}
\begin{equation}
    \mathbf{T}^{n+1}
    =
    \mathbf{T}^{n}
    +
    \frac{\Delta t}{24}
    \left(
        55\mathbf{F}^{n}
        -
        59\mathbf{F}^{n-1}
        +
        37\mathbf{F}^{n-2}
        -
        9\mathbf{F}^{n-3}
    \right).
\end{equation}

Nos métodos de Adams--Bashforth, os primeiros passos foram inicializados com RK4
até que existisse histórico suficiente de derivadas.

\section{Procedimento computacional}

Para cada método, malha e coeficiente $C$, o procedimento numérico foi:
\begin{enumerate}
    \item calcular $\Delta x$, $\Delta y$, $\alpha$ e $\Delta t$ pela
    Equação~\eqref{eq:dt};
    \item inicializar o campo com $T=300\ \si{\kelvin}$ e impor $T=400\ \si{\kelvin}$
    em $x=0$;
    \item calcular $\mathbf{F}(\mathbf{T})$ pelas diferenças finitas da
    Seção 3;
    \item avançar a solução pelo método temporal escolhido;
    \item reaplicar os contornos essenciais em $x=0$ e $x=L_x$;
    \item armazenar, a cada intervalo de aproximadamente $60\ \si{\second}$, as
    temperaturas nos pontos $(L_x/2,L_y/2)$ e $(L_x/4,L_y/4)$;
    \item repetir até $t=4\ \si{\hour}$.
\end{enumerate}

\section{Resultados da Questão 1}

As figuras desta seção mostram a evolução temporal da temperatura nas duas
posições solicitadas, comparando as malhas $10\times10$, $20\times20$,
$40\times40$ e $80\times80$.

\maybeinclude{evolucao_euler_centro_C0.25.png}{Euler explícito em $(L_x/2,L_y/2)$.}{fig:q1_euler_centro}
\maybeinclude{evolucao_euler_quarto_C0.25.png}{Euler explícito em $(L_x/4,L_y/4)$.}{fig:q1_euler_quarto}

\maybeinclude{evolucao_rk2_ponto_medio_centro_C0.25.png}{RK2 ponto médio em $(L_x/2,L_y/2)$.}{fig:q1_rk2pm_centro}
\maybeinclude{evolucao_rk2_ponto_medio_quarto_C0.25.png}{RK2 ponto médio em $(L_x/4,L_y/4)$.}{fig:q1_rk2pm_quarto}

\maybeinclude{evolucao_rk2_euler_modificado_centro_C0.25.png}{RK2 Euler modificado/Ralston em $(L_x/2,L_y/2)$.}{fig:q1_rk2em_centro}
\maybeinclude{evolucao_rk2_euler_modificado_quarto_C0.25.png}{RK2 Euler modificado/Ralston em $(L_x/4,L_y/4)$.}{fig:q1_rk2em_quarto}

\maybeinclude{evolucao_rk2_heun_centro_C0.25.png}{RK2 Heun em $(L_x/2,L_y/2)$.}{fig:q1_heun_centro}
\maybeinclude{evolucao_rk2_heun_quarto_C0.25.png}{RK2 Heun em $(L_x/4,L_y/4)$.}{fig:q1_heun_quarto}

\maybeinclude{evolucao_rk4_centro_C0.25.png}{RK4 clássico em $(L_x/2,L_y/2)$.}{fig:q1_rk4_centro}
\maybeinclude{evolucao_rk4_quarto_C0.25.png}{RK4 clássico em $(L_x/4,L_y/4)$.}{fig:q1_rk4_quarto}

\maybeinclude{evolucao_ab2_centro_C0.12.png}{Adams--Bashforth 2 em $(L_x/2,L_y/2)$.}{fig:q1_ab2_centro}
\maybeinclude{evolucao_ab2_quarto_C0.12.png}{Adams--Bashforth 2 em $(L_x/4,L_y/4)$.}{fig:q1_ab2_quarto}

\maybeinclude{evolucao_ab4_centro_C0.03.png}{Adams--Bashforth 4 em $(L_x/2,L_y/2)$.}{fig:q1_ab4_centro}
\maybeinclude{evolucao_ab4_quarto_C0.03.png}{Adams--Bashforth 4 em $(L_x/4,L_y/4)$.}{fig:q1_ab4_quarto}

\section{Resultados da Questão 2: RK4}

Para RK4, os coeficientes solicitados foram $C=0{,}500$, $C=0{,}250$ e
$C=0{,}125$. Como explicado na Seção 4, o código reescalou esses valores para
$C=0{,}34$, $C=0{,}17$ e $C=0{,}085$ nas execuções finais estáveis.

\maybeinclude{evolucao_rk4_centro_C0.34.png}{RK4 em $(L_x/2,L_y/2)$ para o maior coeficiente estável da série.}{fig:q2_rk4_centro_c034}
\maybeinclude{evolucao_rk4_quarto_C0.34.png}{RK4 em $(L_x/4,L_y/4)$ para o maior coeficiente estável da série.}{fig:q2_rk4_quarto_c034}
\maybeinclude{evolucao_rk4_centro_C0.17.png}{RK4 em $(L_x/2,L_y/2)$ para o coeficiente intermediário.}{fig:q2_rk4_centro_c017}
\maybeinclude{evolucao_rk4_quarto_C0.17.png}{RK4 em $(L_x/4,L_y/4)$ para o coeficiente intermediário.}{fig:q2_rk4_quarto_c017}
\maybeinclude{evolucao_rk4_centro_C0.085.png}{RK4 em $(L_x/2,L_y/2)$ para o menor coeficiente da série.}{fig:q2_rk4_centro_c0085}
\maybeinclude{evolucao_rk4_quarto_C0.085.png}{RK4 em $(L_x/4,L_y/4)$ para o menor coeficiente da série.}{fig:q2_rk4_quarto_c0085}

\maybeinclude{comparacao_C_rk4_centro_Nx20_Ny20.png}{Comparação dos coeficientes $C$ no RK4 para $N_x=N_y=20$, em $(L_x/2,L_y/2)$.}{fig:q2_rk4_comp_centro}
\maybeinclude{comparacao_C_rk4_quarto_Nx20_Ny20.png}{Comparação dos coeficientes $C$ no RK4 para $N_x=N_y=20$, em $(L_x/4,L_y/4)$.}{fig:q2_rk4_comp_quarto}

\section{Resultados da Questão 3: Adams--Bashforth 4}

Para Adams--Bashforth de quarta ordem, os coeficientes solicitados
$C=0{,}500$, $C=0{,}250$ e $C=0{,}125$ foram reescalados para $C=0{,}03$,
$C=0{,}015$ e $C=0{,}0075$.

\maybeinclude{evolucao_ab4_centro_C0.03.png}{AB4 em $(L_x/2,L_y/2)$ para o maior coeficiente estável da série.}{fig:q3_ab4_centro_c003}
\maybeinclude{evolucao_ab4_quarto_C0.03.png}{AB4 em $(L_x/4,L_y/4)$ para o maior coeficiente estável da série.}{fig:q3_ab4_quarto_c003}
\maybeinclude{evolucao_ab4_centro_C0.015.png}{AB4 em $(L_x/2,L_y/2)$ para o coeficiente intermediário.}{fig:q3_ab4_centro_c0015}
\maybeinclude{evolucao_ab4_quarto_C0.015.png}{AB4 em $(L_x/4,L_y/4)$ para o coeficiente intermediário.}{fig:q3_ab4_quarto_c0015}
\maybeinclude{evolucao_ab4_centro_C0.0075.png}{AB4 em $(L_x/2,L_y/2)$ para o menor coeficiente da série.}{fig:q3_ab4_centro_c00075}
\maybeinclude{evolucao_ab4_quarto_C0.0075.png}{AB4 em $(L_x/4,L_y/4)$ para o menor coeficiente da série.}{fig:q3_ab4_quarto_c00075}

\maybeinclude{comparacao_C_ab4_centro_Nx20_Ny20.png}{Comparação dos coeficientes $C$ no AB4 para $N_x=N_y=20$, em $(L_x/2,L_y/2)$.}{fig:q3_ab4_comp_centro}
\maybeinclude{comparacao_C_ab4_quarto_Nx20_Ny20.png}{Comparação dos coeficientes $C$ no AB4 para $N_x=N_y=20$, em $(L_x/4,L_y/4)$.}{fig:q3_ab4_comp_quarto}

\section{Discussão}

Espera-se que o refinamento de malha reduza a diferença entre as curvas, pois a
aproximação espacial usada é de segunda ordem. Como a temperatura no contorno esquerdo foi
mantida em $400\ \si{\kelvin}$ e a condição inicial é $300\ \si{\kelvin}$, as
temperaturas internas inicialmente tendem a aumentar devido ao aquecimento
imposto em $x=0$. A evolução final é afetada também pelas condições convectivas
e radiativas em $y=0$ e $y=L_y$.

Os métodos RK2 e RK4 usam mais avaliações de $\mathbf{F}$ por passo do que Euler,
mas reduzem o erro temporal para um mesmo $\Delta t$. Os métodos
Adams--Bashforth reaproveitam avaliações anteriores da derivada, tendo custo
baixo por passo depois da inicialização; por outro lado, exigem histórico e uma
escolha mais conservadora de $C$, especialmente no caso AB4.

\section{Conclusão}

Foram implementados os sete métodos explícitos pedidos e aplicados ao problema
de condução bidimensional transiente. A mesma discretização espacial por
diferenças finitas foi mantida em todos os casos, permitindo comparar apenas o
efeito do método temporal, da malha e do coeficiente de estabilidade.

Os arquivos CSV e PNG produzidos por \texttt{scripts/run\_all.py} documentam as
séries temporais em $(L_x/2,L_y/2)$ e $(L_x/4,L_y/4)$ para todas as combinações
pedidas no enunciado.

\section{Referências}

\begin{thebibliography}{9}
\bibitem{chapra}
S. C. Chapra e R. P. Canale.
\textit{Numerical Methods for Engineers}. McGraw-Hill.

\bibitem{incropera}
F. P. Incropera, D. P. DeWitt, T. L. Bergman e A. S. Lavine.
\textit{Fundamentals of Heat and Mass Transfer}. Wiley.

\bibitem{burden}
R. L. Burden e J. D. Faires.
\textit{Numerical Analysis}. Cengage Learning.
\end{thebibliography}

\end{document}



